\documentclass[11pt,a4paper]{article}

\usepackage[T1]{fontenc}
\usepackage[utf8]{inputenc}
\usepackage[english]{babel}
\usepackage{lmodern}
\usepackage{microtype}
\usepackage[margin=2.5cm]{geometry}
\usepackage{booktabs}
\usepackage{array}
\usepackage{enumitem}
\usepackage{xcolor}
\usepackage{listings}
\usepackage{hyperref}
\usepackage{url}

\hypersetup{
  colorlinks=true,
  linkcolor=blue!55!black,
  urlcolor=blue!55!black,
  citecolor=blue!55!black,
  pdftitle={Persistent Causal State for Operating-System Enforcement}
}

\lstdefinelanguage{json}{
  basicstyle=\ttfamily\small,
  string=[s]{"}{"},
  stringstyle=\color{green!35!black},
  comment=[l]{//},
  commentstyle=\color{gray},
  breaklines=true,
  columns=fullflexible,
  keepspaces=true,
  showstringspaces=false
}

\title{Persistent Causal State for Operating-System Enforcement\\
\large Engineering a Stateful, Auditable, and Fail-Closed Linux Policy Runtime}

\author{
Felipe Maya Muniz\\
\textit{Unix-AGB Project / AletheionAGI}
}

\date{August 17, 2026}

\begin{document}
\maketitle

\begin{abstract}
Most operating-system enforcement mechanisms decide from the current request:
an executable, a credential, a path, a socket address, or a static label. This
local view is efficient and dependable, but it cannot directly express risk
that emerges from a sequence of individually ordinary actions. Unix-AGB
(Aletheion Guard Bridge) investigates a complementary architecture in which
normalized kernel observations update persistent, namespace-isolated causal
state; a deterministic policy layer then compiles only restrictive decisions
into short-lived enforcement records. Exact events remain the historical
truth, learned state remains replaceable, and Linux retains final enforcement
authority. The key design constraint is \emph{monotonic authority}: learned or
stateful inference may only reduce existing privilege; it cannot manufacture
permission. This article explains the engineering of that pipeline: contracts,
identity, ordering, persistence, policy revisions, authenticated caches,
seccomp user-notification handoff, recovery, packaging, and experimental
validation. It also states the boundary of the result. The current system is a
controlled prototype with packaged, opt-in, exact-launch egress enforcement;
it is not a production security service and does not yet establish security
efficacy on natural unknown attacks.
\end{abstract}

\noindent\textbf{Keywords:}
persistent causal state,
operating-system security,
stateful enforcement,
monotonic authority,
deny-only compilation,
Linux security,
seccomp,
runtime security.

\section{The engineering problem}

Consider two processes that perform the same terminal operation: opening a
credential file or connecting to an external address. One process may have
followed a routine configuration path; the other may have executed an unusual
tool, established external connectivity, and only then accessed the sensitive
resource. A stateless rule sees identical terminal requests. A stateful system
can distinguish the histories, but adding state to operating-system policy
creates difficult requirements:

\begin{itemize}[leftmargin=*]
  \item events must be ordered and attributed to a stable identity;
  \item state must survive restart without accepting corruption or rollback;
  \item model output must not become historical truth or silently grant access;
  \item enforcement must remain deterministic, bounded, and available during
        control-plane failures;
  \item every negative action must be linked to durable evidence and an exact
        policy revision;
  \item recovery must not broaden a denial to unrelated processes.
\end{itemize}

Unix-AGB approaches these constraints by separating observation, canonical
evidence, state inference, policy compilation, and kernel enforcement. The
separation is more important than any individual state engine.

In this paper, \emph{causal state} denotes trajectory-dependent state whose
policy interpretation depends on ordered predecessor events within the same
security namespace; we do not claim causal identification in the statistical
or interventionist sense. Unix-AGB is therefore not primarily a new trajectory
detector. It is an architecture for safely compiling persistent,
namespace-scoped trajectory state into short-lived kernel-facing restrictions.
The state engine determines when a trajectory warrants restriction; Unix-AGB
determines how that conclusion can affect kernel enforcement without giving the
state engine authority to expand privilege.

\subsection*{Contributions}

This work makes four contributions:

\begin{enumerate}[leftmargin=*]
  \item an operating-system policy architecture that separates persistent
        causal state from canonical evidence and deterministic enforcement;
  \item deny-only compilation governed by the Monotonic Authority Principle,
        under which learned state may restrict but cannot expand privilege;
  \item a restart- and revision-aware enforcement pipeline with authenticated
        short-lived policy records and durable negative audit; and
  \item an experimental Linux prototype demonstrating namespace-scoped seccomp
        enforcement, recovery, rollback, and failure semantics.
\end{enumerate}

\section{Related Work and Positioning}

Linux Security Modules provide a common kernel framework for mandatory access
control systems such as SELinux and AppArmor~\cite{lsm,selinux,apparmor}.
Seccomp user notification permits selected system calls to be mediated in user
space~\cite{seccomp}, while eBPF-based systems move programmable observation
and runtime detection close to kernel events~\cite{falco}. Unix-AGB does not
replace these mechanisms: it uses existing Linux enforcement surfaces as the
last step of a separately versioned state-and-policy pipeline.

Prior work has reconstructed attack histories from operating-system
provenance~\cite{backtracker,camflow}, tracked information flow and
taint~\cite{denning}, correlated event streams with stateful and complex-event
processing~\cite{cep}, and learned anomalous system-call behavior
~\cite{forrest}. Commercial behavior-based endpoint detection follows a
related operational pattern, although its internal semantics are usually not a
public systems contract. These approaches establish that histories can improve
detection and investigation. The distinct question addressed here is how a
persistent trajectory conclusion becomes a compact, authenticated, expiring,
revision-bound kernel restriction that cannot grant privilege and remains
auditable across restart, rollback, and component failure. Detection quality is
deliberately outside the central claim: FSMs, sequence rules, windows, proxies,
and ASM-CM all enter through the same state-engine boundary.

ActPlane is particularly close in motivation: it bridges agent-level policy
context to OS-level enforcement, uses an information-flow DSL for cross-event
policies, and implements kernel mediation with eBPF~\cite{actplane}. Unix-AGB
addresses a complementary systems boundary. Rather than allowing an agent to
declare the kernel policy, it treats trajectory state as derived and
potentially learned input to a deterministic compiler. The resulting records
are deny-only, authenticated, revision-bound, expiring, and tied to durable
canonical evidence. Its emphasis is therefore persistent state across restart
and policy revision, monotonic authority, auditable negative decisions, and
explicit recovery and failure semantics.

\section{System decomposition}

The architecture is a pipeline of narrow trust boundaries:

\begin{center}
\begin{tabular}{>{\raggedright\arraybackslash}p{0.20\textwidth} >{\raggedright\arraybackslash}p{0.70\textwidth}}
\toprule
Layer & Responsibility \\
\midrule
Linux and collectors & Observe process, file, identity, and network operations; provide stable kernel facts and enforcement hooks. \\
Event gateway & Validate schema, identity, ordering, provenance, and size before accepting an event. \\
Canonical store & Preserve exact append-only facts. Reject duplicates and sequence replay. \\
State engine & Maintain bounded state per security namespace and emit a summary with revision, signals, evidence identifiers, and optional confidence. \\
Policy engine & Apply deterministic invariants and translate validated state into \texttt{ALLOW}, \texttt{DENY}, or \texttt{ABSTAIN}. \\
Compiled cache & Retain only authenticated, versioned, expiring \texttt{DENY} records for the enforcement edge. \\
Enforcement adapter & Bind a compiled decision to an exact kernel-visible subject and return a concrete result such as \texttt{EACCES}. \\
Optional reader & Produce human-readable explanations outside the synchronous enforcement path. \\
\bottomrule
\end{tabular}
\end{center}

The canonical store and the state engine deliberately answer different
questions. The store answers ``what exactly was observed?'' The state engine
answers ``what bounded security state follows from the accepted history?'' A
checkpoint, neural representation, or rule-engine accumulator is therefore a
derived artifact. It may be discarded and reconstructed; canonical evidence
must not be silently rewritten to agree with it.

\section{Canonical events and stable identity}

\subsection{The event contract}

Each accepted event carries a schema version, globally unique event identifier,
per-namespace sequence, wall and monotonic timestamps, host identifier,
namespace, subject, operation, resource, observed or requested result, policy
revision, labels, and provenance. The first operation vocabulary contains
\texttt{process.exec}, \texttt{process.exit}, \texttt{file.open},
\texttt{network.connect}, and \texttt{identity.change}.

The contract is intentionally metadata-oriented. It does not require file
contents, packet payloads, secret values, or prompt bodies. Content resolution
would be a separate authorized operation. This reduces both privacy exposure
and the amount of attacker-controlled material admitted to the policy plane.

\subsection{PID reuse is an identity bug}

A numeric process identifier is recyclable. State keyed only by PID can leak a
previous process's history into a new process. Unix-AGB instead uses:

\begin{center}
\texttt{process:<boot-id>:<pid>:<start-time-ns>}
\end{center}

The boot identifier separates reboots; start time separates PID reuse inside a
boot. The gateway recomputes the expected namespace from the subject fields and
rejects mismatches. In the Gate 4 handoff, the launcher and guardian derive the
identity independently from \texttt{/proc}; a handoff is rejected if the
claimed namespace does not match the live process.

\subsection{Ordering and gaps}

Sequence is strictly increasing within a namespace. Duplicate event IDs and
non-increasing sequences are rejected before the in-memory index changes. A
gap is not interpreted as benign absence. It marks the state incomplete until
an explicit resynchronization policy is applied. Wall-clock timestamps are
descriptive; sequence and monotonic time define causal order.

This rule prevents a common stateful-security failure: treating lost telemetry
as evidence that nothing happened.

\section{Persistent causal state}

\subsection{Namespace-isolated transition state}

The state engine consumes one validated event at a time and emits a
\texttt{SecurityStateSummary}. Its revision must match the accepted event
sequence. A summary contains a risk band, signals, evidence IDs, engine name,
checkpoint fingerprint when applicable, and update time. Queries and updates
require an exact namespace match; state is never implicitly shared across
processes, users, containers, services, agents, or tenants.

A minimal deterministic example can remember that an execution occurred,
then that an external connection followed it, and finally that a credential was
accessed. The terminal credential event becomes elevated only when the earlier
relations exist in the same namespace. The important property is not the
specific rule, but the state transition contract:

\[
S_{n+1} = F(S_n, E_{n+1}, C, R)
\]

where $S_n$ is namespace state, $E_{n+1}$ is the next canonical event, $C$ is
the configuration fingerprint, and $R$ is the engine or checkpoint revision.
A consumer must be able to identify all four.

\subsection{State engines are replaceable}

Unix-AGB exposes a state-engine boundary rather than coupling policy to one
model. Deterministic finite-state, sequence-rule, sliding-window, stateful
proxy, and ASM-CM adapters can be evaluated through the same event and summary
contracts. This permits baseline comparisons and prevents learned inference
from redefining the policy protocol.

The research history matters. Early causal-generalization experiments produced
negative results; canonical trajectory-local identifiers later resolved a
representation failure, and a frozen three-seed ensemble succeeded on a
synthetic held-out test. Those results support a mechanism and a controlled
family, not detection of natural unknown attacks. The operational design
therefore defaults model disagreement to \texttt{ABSTAIN} and keeps neural
inference outside the syscall hot path.

\subsection{Atomic snapshots and restore}

Persistent state is useful only if restore is stricter than ordinary loading.
The prototype snapshot protocol binds:

\begin{itemize}[leftmargin=*]
  \item snapshot format version;
  \item engine fingerprint;
  \item configuration fingerprint;
  \item complete, deterministically ordered namespace state;
  \item checksum or authentication data.
\end{itemize}

Checkpoint writes use a temporary file in the destination directory, flush and
\texttt{fsync}, atomic rename, and directory \texttt{fsync}. State updates are
rolled back in memory if checkpoint persistence fails. Restore rejects corrupt,
unsupported, wrong-engine, or wrong-configuration snapshots rather than
inventing permissive state.

Checksums detect accidental corruption. Security-sensitive policy snapshots
add HMAC authentication because an attacker able to rewrite state can also
recompute an ordinary hash.

\section{From state to policy}

\subsection{A small enforcement algebra}

The policy core deliberately uses only three effects:

\begin{description}[leftmargin=2.4cm,style=nextline]
  \item[\texttt{ALLOW}] No restricting invariant was found. In the current
        Gate 3 design this is a shadow/audit result, not a compiled grant.
  \item[\texttt{DENY}] Validated state and policy justify a restriction.
  \item[\texttt{ABSTAIN}] Evidence, confidence, revision, persistence, or
        contract integrity is insufficient for a trustworthy decision.
\end{description}

Warnings are audit metadata, not a fourth kernel action. This keeps the policy
algebra small and ensures domain-specific assessments are translated by an
adapter before reaching the enforcement layer.

\subsection{Decision order and monotonic authority}

Gate 3 first validates policy configuration, event and state contracts,
namespace equality, policy revision, and state revision. Static
\texttt{restricted} or \texttt{quarantined} state yields \texttt{DENY}.
Elevated learned state requires evidence, minimum confidence, and a checkpoint
fingerprint. Normal or monitor state yields a shadow \texttt{ALLOW}. Unknown,
stale, incomplete, or inconsistent input yields fail-closed
\texttt{ABSTAIN}.

Only \texttt{DENY} is compiled. Neither learned state nor an unavailable cache
can manufacture permission.

\paragraph{Monotonic Authority Principle.}
State-derived decisions may preserve or reduce authority granted by the base
policy, but never expand it. Base Linux policy therefore remains authoritative
when inference abstains, state expires, or the stateful control plane becomes
unavailable. Deny-only compilation is the concrete enforcement rule that
realizes this principle in Unix-AGB.

\subsection{Durability before restriction}

A negative decision is appended to the audit log and synchronized before it
becomes cacheable. If audit persistence fails, the result becomes
\texttt{ABSTAIN} with an explicit reason code. Audit-only positive and abstain
records may be group-synchronized for throughput, but a \texttt{DENY} forces
durability immediately.

This ordering provides a recoverable invariant:

\begin{center}
\textit{if enforcement can observe a compiled denial, durable audit evidence
for that denial already exists.}
\end{center}

\section{Authenticated decision compilation}

A compiled decision binds the full namespace, operation, SHA-256 digest of the
canonical resource, policy revision, state revision, evidence digest, decision
identifier, and expiry. The cache accepts only \texttt{DENY}. Its snapshot
authenticates the ordered record set and active policy revision with
HMAC-SHA-256.

Cache lookup returns \texttt{ABSTAIN} on a miss, expiry, policy mismatch, or
stale state revision. Rotation and rollback clear incompatible in-memory state.
An older snapshot cannot be relabeled as a newer policy. The Gate 4 runtime
projection additionally retains the last authenticated snapshot when a reload
is malformed; initial invalid policy prevents activation, while an authenticated
empty rotation or expiry removes the restriction.

The current service-supervision prototype maps any active trajectory-level
\texttt{DENY} to external-egress containment for the exact process namespace.
The denial may originate at a credential read, persistence write,
administrative execution, or network event; its originating operation remains
in the audit. This mapping is intentionally broad and must not be called
destination-specific: the runtime checks namespace, revision, and expiry but
does not compare the intercepted destination with the compiled resource digest.

\section{The Gate 4 enforcement edge}

\subsection{Install interception before execution}

The opt-in launcher creates the protected child and installs a seccomp filter
with user notification before releasing it into \texttt{exec}. This ordering
closes a race in which the application could connect before enforcement was
attached. The listener file descriptor and a bounded memory-reading adapter are
transferred to a long-lived guardian over a Unix-domain control socket.

The handoff is authenticated with peer credentials, an HMAC, a policy revision,
a nonce, and a short monotonic expiry. File descriptors are passed with
\texttt{SCM\_RIGHTS}. The guardian also verifies process-group scope and
independently recomputes the target namespace. Replay, wrong peer, stale
revision, expired message, namespace mismatch, or incomplete descriptor transfer
causes rejection.

\subsection{Decision at \texttt{connect}}

For each seccomp notification the guardian:

\begin{enumerate}[leftmargin=*]
  \item validates that the notification is still live;
  \item resolves the notifying thread to its thread-group process;
  \item checks exact target scope;
  \item reads at most 128 bytes of the requested \texttt{sockaddr};
  \item permits loopback and Unix-domain destinations according to the frozen
        policy;
  \item consults the current authenticated Gate 3 cache for external egress;
  \item appends and synchronizes the enforcement audit record;
  \item returns \texttt{EACCES} for a matched denial or continues the syscall
        otherwise.
\end{enumerate}

Unknown address families, malformed lengths, failed memory reads, or adapter
failure take a scoped fail-closed path for the protected target. Out-of-scope
processes do not inherit the target's denial.

\subsection{Listener ownership and recovery}

Early experiments exposed a critical distinction between a policy-worker crash
and listener loss. If the process owning a seccomp notification listener dies,
the protected syscall does not necessarily receive a clean \texttt{EACCES}; it
may remain blocked or fail through kernel teardown semantics. Calling that
behavior ``fail closed'' would hide an availability failure.

The engineering response was to retain listener ownership in a guardian while
making the policy worker replaceable. In-flight notification IDs are validated
before a response; worker failures consume a bounded restart budget and use a
scoped fallback. Pre-lease deadlines and scoped process teardown bound cases in
which no worker has safely leased the notification. Recovery never broadens a
target denial to an unrelated process.

\section{Packaging and operation}

The laboratory runtime is packaged as an opt-in Debian artifact with a systemd
unit, dedicated system account, configuration under \texttt{/etc/unix-agb},
runtime sockets under \texttt{/run/unix-agb}, state and audit under
\texttt{/var/lib/unix-agb}, and root-owned authentication keys. Configuration
is revision-bound and disabled by default. The service applies systemd
hardening such as \texttt{NoNewPrivileges}, \texttt{ProtectSystem=strict},
private temporary space, capability bounding, and explicit writable paths.

Package lifecycle tests cover install, inactive default, explicit activation,
restart, daemon reload, reboot, upgrade, disable, purge, dedicated-account
cleanup, socket cleanup, and residue auditing. A launcher readiness race was
found during negative testing: systemd reported a simple service active before
its control socket existed. The launcher now waits for a bounded interval and
fails before releasing the protected child if readiness is not reached.

This is a useful example of why packaging is part of security engineering.
Correct in-process logic does not establish safe boot ordering, upgrade
behavior, secret ownership, or recovery.

\section{Experimental evidence}

The repository preserves protocols, failed attempts, environment fingerprints,
and reports rather than reporting only successful demonstrations. The relevant
engineering progression includes:

\begin{itemize}[leftmargin=*]
  \item a causal proof with identical terminal operations and divergent
        trajectory-dependent shadow decisions;
  \item dry-run Gate 3 policy with revision checks, durable audit, authenticated
        deny-only cache, expiry, corruption rejection, and namespace isolation;
  \item a seccomp user-notification pilot returning \texttt{EACCES};
  \item concurrent protected calls, overload tests, timeout paths, and the
        negative listener-loss result;
  \item retained-listener guardian and replaceable worker recovery;
  \item packaged, opt-in exact-launch enforcement on controlled Ubuntu 24.04
        and 26.04 virtual machines;
  \item a controlled long-lived service transition from authenticated empty
        policy, to exact-namespace denial, through corrupt reload retention,
        and back to authenticated empty policy.
\end{itemize}

Measured samples range from tens of microseconds in the retained-listener
guardian experiments to low milliseconds in concurrent broker paths. These
numbers characterize particular controlled boundaries; they are not production
capacity claims. The service-supervision run observed the protected application
change from an allowed external attempt to \texttt{EACCES}, retain denial after
a malformed cache replacement, and recover after authenticated empty rotation,
while an unprotected control remained unaffected.

\begin{center}
\small
\begin{tabular}{>{\raggedright\arraybackslash}p{0.43\textwidth} >{\raggedright\arraybackslash}p{0.47\textwidth}}
\toprule
Experiment & Observed result \\
\midrule
Retained-listener guardian, 256 target calls~\cite{guardian} & 256/256 returned \texttt{EACCES}; decision latency 44~$\mu$s p50, 49~$\mu$s p95, 68~$\mu$s p99 \\
Concurrent broker, 256 target calls~\cite{concurrency} & 256/256 returned \texttt{EACCES}; 1.148~ms p50, 1.939~ms p95, 2.366~ms p99 \\
In-flight worker failure~\cite{inflight} & Exact pending notification returned \texttt{EACCES} in 123~$\mu$s; replacement worker ready in 9.839~ms \\
Exact-launch VM enforcement~\cite{exactlaunch} & Protected external connect returned \texttt{EACCES}; loopback and unprotected control remained unaffected \\
Authenticated cache transitions & Corrupt reload retained the last valid denial; authenticated empty rotation removed it \\
Service-supervision audit & 191~$\mu$s p50 and 276~$\mu$s p95/p99, including durable audit \\
\bottomrule
\end{tabular}
\end{center}

The table reports separate experimental boundaries and must not be read as one
end-to-end latency distribution. In particular, the guardian microbenchmark,
concurrent broker, packaged VM, and service-supervision runs used different
controlled workloads.

A subsequent transition replaced the handcrafted trigger with the frozen
three-member ensemble and the real Gate 3 policy compiler while the protected
service remained alive. The resulting authenticated \texttt{DENY}, originating
at \texttt{file.open}, caused the next external connect to receive
\texttt{EACCES}. Its input was nevertheless a frozen controlled trajectory
explicitly rebound to the live namespace, not newly captured BPF telemetry.
The frozen promotion matrix therefore continues to classify Gate 4 as a
controlled prototype.

\section{Failure semantics}

``Fail closed'' must name both scope and consequence. Unix-AGB distinguishes:

\begin{itemize}[leftmargin=*]
  \item \textbf{policy abstention}: no trustworthy state-derived decision;
  \item \textbf{scoped denial}: the protected request receives
        \texttt{EACCES};
  \item \textbf{scoped teardown}: the protected workload is terminated when a
        safe response cannot be guaranteed;
  \item \textbf{base-policy continuity}: Linux policy remains active if the
        model or control plane fails;
  \item \textbf{availability failure}: the request hangs or loses a usable
        response, which must not be mislabeled as successful denial.
\end{itemize}

This vocabulary forces the design to account for missing telemetry, sequence
gaps, corrupt snapshots, expired decisions, worker crashes, listener loss,
audit failure, policy rollback, and restart races independently.

\section{Security properties and residual risks}

The implemented contracts support several concrete properties:

\begin{itemize}[leftmargin=*]
  \item stable process identity resists PID reuse;
  \item exact namespace checks limit cross-process state leakage;
  \item replay and sequence checks preserve event order;
  \item authenticated, revision-bound cache snapshots resist unauthenticated
        policy replacement;
  \item deny-only compilation prevents model-derived privilege expansion;
  \item expiry and authenticated empty rotation provide bounded rollback;
  \item durable audit precedes compiled negative enforcement;
  \item independent handoff verification binds the listener to the live target.
\end{itemize}

Residual risks remain substantial. A compromised privileged collector can
forge observations; a compromised kernel can defeat both observation and
enforcement; HMAC keys require production-grade provisioning and rotation;
destination-specific binding is incomplete in the current service projection;
system-wide non-cooperative attachment is not implemented; natural-telemetry
efficacy remains unproven; and broad fail-closed policy can become a denial-of-
service mechanism. These are promotion blockers, not documentation details.

\section{Why persistent causal state belongs beside, not inside, the kernel}

The prototype does not yet justify a custom kernel. Linux already provides the
necessary experimental surfaces: BPF or audit-style observation, stable process
metadata, Unix peer credentials, cgroups, systemd supervision, seccomp user
notification, and established LSMs. Causal state, checkpoint management,
policy revisioning, review, and optional learned inference evolve faster and
have larger failure surfaces than a minimal kernel hook.

Keeping that machinery in user space permits restart, replacement, comparison
with deterministic baselines, and explicit abstention. The kernel-facing edge
receives a compact, expiring, authenticated restriction rather than a model.
Kernel work should be considered only if measurement identifies a missing hook,
unavailable identity, unacceptable latency, or an enforcement property that
existing Linux interfaces cannot provide.

\section{Conclusion}

Persistent causal state can enrich operating-system enforcement only when it is
engineered as a constrained evidence and policy pipeline, not as an opaque model
placed in the syscall path. Unix-AGB's central pattern is:

\begin{center}
exact events $\rightarrow$ isolated persistent state $\rightarrow$ validated
summary $\rightarrow$ deterministic policy $\rightarrow$ authenticated,
expiring denial $\rightarrow$ scoped kernel enforcement.
\end{center}

The durable boundaries are the contribution: canonical facts remain separate
from derived state; identity survives PID reuse; uncertainty becomes
\texttt{ABSTAIN}; only denial is compiled; audit precedes restriction; and the
listener remains available while replaceable policy machinery recovers. The
combination of Persistent Causal State, Deny-Only Compilation, and Monotonic
Authority is intended to apply beyond this particular runtime. The current
evidence shows that these mechanisms can be assembled and tested in a
packaged, opt-in Linux prototype. It does not yet show production readiness or
general unknown-attack detection. Advancing from mechanism to security claim
requires the frozen promotion matrix: independent live Gate 3 decision binding,
concurrent namespace isolation, endurance and resource bounds, authenticated
update and rollback, failure injection, real application pilots, and multi-VM
lifecycle validation under one fixed revision and artifact.

\begin{thebibliography}{99}

\bibitem{lsm}
C. Wright et al.,
\emph{Linux Security Modules: General Security Support for the Linux Kernel},
11th USENIX Security Symposium, 2002.

\bibitem{selinux}
P. Loscocco and S. Smalley,
\emph{Integrating Flexible Support for Security Policies into the Linux
Operating System}, USENIX Annual Technical Conference, 2001.

\bibitem{apparmor}
Linux Kernel Documentation,
\emph{AppArmor}, \url{https://docs.kernel.org/admin-guide/LSM/apparmor.html}.

\bibitem{seccomp}
Linux Kernel Documentation,
\emph{Seccomp BPF User Notification},
\url{https://docs.kernel.org/userspace-api/seccomp_filter.html}.

\bibitem{falco}
The Falco Authors,
\emph{Falco: Cloud Native Runtime Security}, \url{https://falco.org/docs/}.

\bibitem{backtracker}
S. T. King and P. M. Chen,
\emph{Backtracking Intrusions}, 19th ACM Symposium on Operating Systems
Principles, 2003.

\bibitem{camflow}
T. Pasquier et al.,
\emph{Practical Whole-System Provenance Capture}, ACM Symposium on Cloud
Computing, 2017.

\bibitem{denning}
D. E. Denning,
\emph{A Lattice Model of Secure Information Flow}, Communications of the ACM,
19(5), 1976.

\bibitem{cep}
G. Cugola and A. Margara,
\emph{Processing Flows of Information: From Data Stream to Complex Event
Processing}, ACM Computing Surveys, 44(3), 2012.

\bibitem{forrest}
S. Forrest et al.,
\emph{A Sense of Self for Unix Processes}, IEEE Symposium on Security and
Privacy, 1996.

\bibitem{actplane}
Y. Zheng et al.,
\emph{ActPlane: Programmable OS-Level Policy Enforcement for Agent Harnesses},
arXiv:2606.25189v2, 2026,
\url{https://arxiv.org/abs/2606.25189v2}.

\bibitem{architecture}
Unix-AGB Project,
\emph{Unix-AGB Architecture Specification},
\texttt{docs/Unix\_AGB\_Architecture\_Specification\_PTBR.md}, 2026.

\bibitem{contracts}
Unix-AGB Project,
\emph{AGB Contracts},
\texttt{docs/contracts.md}, 2026.

\bibitem{gate3}
Unix-AGB Project,
\emph{Gate 3 Dry-Run Policy Engine},
\texttt{docs/gate3-policy-engine.md}, 2026.

\bibitem{promotion}
Unix-AGB Project,
\emph{Gate 4 Promotion Protocol},
\texttt{docs/report/106\_gate4\_promotion\_protocol\_20260817.md}, 2026.

\bibitem{supervision}
Unix-AGB Project,
\emph{Gate 4 Gate 3 Service Supervision},
\texttt{docs/report/108\_gate4\_gate3\_service\_supervision\_20260817.md}, 2026.

\bibitem{concurrency}
Unix-AGB Project,
\emph{Gate 4 Concurrent Seccomp Broker},
\texttt{docs/report/94\_gate4\_concurrent\_seccomp\_broker\_20260817.md}, 2026.

\bibitem{guardian}
Unix-AGB Project,
\emph{Gate 4 Retained-Listener Guardian},
\texttt{docs/report/96\_gate4\_retained\_listener\_guardian\_20260817.md}, 2026.

\bibitem{inflight}
Unix-AGB Project,
\emph{Gate 4 In-Flight Notification Recovery},
\texttt{docs/report/97\_gate4\_inflight\_notification\_recovery\_20260817.md}, 2026.

\bibitem{exactlaunch}
Unix-AGB Project,
\emph{Gate 4 Persistent Exact-Launch Enforcement},
\texttt{docs/report/105\_gate4\_persistent\_exact\_launch\_enforcement\_20260817.md}, 2026.

\end{thebibliography}

\end{document}
